\documentclass[mode=chapternotes, palette=purple, fontsize=12pt, font=default, official=false]{modernclassnotes}

\course{Boolean Algebra}
\coursecode{MATH-405}
\professor{Leonhard Euler}
\institution{Department of Mathematics}
\term{Fall 2026}
\docsubtitle{Chapter Notes}
\title{Notes on Required Readings}

\begin{document}

\maketitle
\tableofcontents

\newpage
% ===== Chapter 1 =====
\chapter{1}{Digital Systems and Binary Numbers}{%
  Understand the binary number system.,%
  {Know how to convert between binary, octal, decimal, and hexadecimal numbers.},%
  Know how to take the complement and reduced radix complement of a number.,%
  Know how to form the code of a number.,%
  Know how to form the parity bit of a word.%
}
% ----- Section 1 -----
\begin{itemize}
    \item The \textbf{binary system} has just two \textbf{digits (bits)}, 0 and 1. Discrete elements of information can be represented with group of these bits called \textbf{binary codes}
    \item You can indicate the base when clarification is needed, e.g. \(0111_{2}=7_{10}\)
    \item \textit{Modern digital design methodology} uses \textbf{hardware description languages (HDLs)} to describe and simulate digital circuits in a textual form
    \item \textbf{Digital (logic) circuits} process \textit{binary data} by using \textbf{\textit{binary logic elements (logic gates)}}
    \begin{itemize}
        \item Quantities are stored in binary storage elements \textbf{(flip-flops)}
    \end{itemize}
\end{itemize}
% ----- Section 2 -----
\subchapter{2}{Binary Numbers}

% ----- Section 3 -----
\subchapter{3}{Number-Base Conversions}

% ----- Section 4 -----
\subchapter{4}{Octal and Hexadecimal Numbers}

% ----- Section 5 -----
\subchapter{5}{Complements of Numbers}

% ----- Section 6 -----
\subchapter{6}{Signed Binary Numbers}

% ----- Section 7 -----
\subchapter{7}{Binary Codes}

% ----- Section 8 -----
\subchapter{8}{Binary Storage and Registers}

% ----- Section 9 -----
\subchapter{9}{Binary Logic}


\newpage
% ===== Chapter 2 =====
\chapter{2}{Boolean Algebra and Logic Gates}{%
  Gain a basic understanding of postulates used to form algebraic structures.,%
  Understand the Huntington Postulates.,%
  Understand the basic theorems and postulates of Boolean algebra.,%
  Know how to develop a logic diagram from a Boolean function; know how to derive a Boolean function from a logic diagram.,%
  Know how to apply DeMorgan's theorems.,%
  Know how to express a Boolean function as a truth table; know how to derive a Boolean function from a truth table.,%
  Know how to express a Boolean function as a sum of minterms and as a product of maxterms.,%
  {Be able to convert from a sum of minterms to a product of maxterms, and vice versa.},%
  Be able to form a two-level gate structure from a Boolean function in sum of products form; know how to form a two-level gate structure from a Boolean function in product of sums form.,%
  Be able to implement a Boolean function with NAND and inverter gates; know how to implement a Boolean function with NOR and inverter gates.%
}
% ----- Section 1 -----
\subchapter{1}{Introduction}
\begin{itemize}
\item Placeholder \( x = \frac{1}{2}\)
\end{itemize}

\end{document}
